\section{System and Experimental Protocol}
\label{sec:method}

\subsection{Modernized Pluggable System}

We modernize the reference DeepSORT implementation \cite{bewley2016sort,wojke2017deepsort}
along exactly two axes -- detection and appearance -- while leaving the tracking core untouched.
The original codebase consumes \emph{precomputed} detections from disk and a fixed appearance
encoder. We replace both with pluggable, swappable adapters selected by name at run time: a
detector front-end and a re-identification (REID) front-end, each registered through a small
registry so a tracker is specified as a (detector, REID) pair. In this paper the detector is
YOLOv8m \cite{jocher2023yolov8} and the REID model is either OSNet \cite{zhou2019osnet} (trained
for person re-identification) or a generic ImageNet \texttt{timm} backbone \cite{wightman2019timm}
(\emph{not} trained for re-identification), used as a deliberately weak appearance baseline.

The SORT core is preserved verbatim: an 8-dimensional constant-velocity Kalman filter, Hungarian
assignment, and a cascaded matcher that combines a Mahalanobis motion gate with a cosine
\emph{appearance} gate over REID descriptors. Its track lifecycle (confirmation after $n_{\text{init}}$
hits, deletion after $\text{max\_age}$ misses) is unchanged. The integration contract between the
modern front-ends and this core is a single value object: each detection is a
\texttt{Detection(tlwh, confidence, feature)} -- a top-left/width-height box, a scalar detector
confidence, and an appearance vector -- and the per-frame loop is the original
\texttt{tracker.predict()} followed by \texttt{tracker.update(detections)}. Because the contract and
the core are fixed, any change in tracking quality is attributable to the detector, the REID model,
or the pre-tracker \emph{gating} that decides which detections reach \texttt{update()}.

\paragraph{The dropped gate.}
This last point is central. In the original pipeline, precomputed detections passed through a fixed
gate before reaching the tracker~\cite{wojke2017deepsortcode}: a confidence threshold
($\text{min\_confidence}$), a non-maximum-suppression (NMS) step on overlapping boxes, and a minimum
box-height filter. The
modern live pipeline runs the detector \emph{inside} the loop, and its default configuration
\textbf{drops this gate}: $\text{min\_confidence}=0.0$ and $\text{nms\_max\_overlap}=1.0$, which
disables both the confidence filter and the secondary NMS (on the assumption that the detector's own
internal NMS and \texttt{conf} threshold suffice). Every YOLO detection above the detector's own
low internal threshold is therefore lifted directly into a \texttt{Detection} and offered to the
tracker. This single configuration change -- not any change to the Kalman filter or the matcher -- is
the regression we audit in \cref{sec:results}: ungated low-confidence boxes are free to instantiate
spurious new tracks. The gate itself is still present in the code path (the same confidence-filter-plus-NMS
gate as the reference pipeline~\cite{wojke2017deepsortcode}), merely disabled by default, which makes the
restoration ablation below a one-line configuration change rather than a code change.

\subsection{Evaluation Protocol}

We evaluate on six MOT-Challenge training sequences spanning two benchmarks: TUD-Campus,
TUD-Stadtmitte, KITTI-17, and PETS09-S2L1 from 2D MOT 2015 \cite{leal2015mot15}, and MOT16-09 and
MOT16-11 from MOT16 \cite{milan2016mot16}. The set deliberately mixes clean, sparse scenes with
dense, low-resolution ones, since per-clip density is the axis along which our central effect
reverses sign. The primary metric is HOTA \cite{luiten2021hota}, computed by the official TrackEval
toolkit \cite{luiten2020trackeval}; we report its DetA and AssA sub-scores when decomposing detection
versus association. End-to-end throughput, timed on one sequence (MOT16-09), is \fpsRealtime{}~FPS
(detection and re-identification each run near $30$~FPS; their serial composition is the bottleneck).

\subsection{Controlled Experimental Design}

The experiments are designed so that each varies one lever while holding the others fixed.

\paragraph{(a) Detection-vs-association decomposition (GT-box, REID-only).}
To separate the modernization gain into a detection component and an association component, we run a
\emph{ground-truth-box} mode in which the detector is replaced by an oracle that emits the ground-truth
boxes, with the confidence/NMS gate skipped. Detection is then perfect and identical across runs, so
the only free variable is the REID model feeding the appearance gate. Comparing OSNet against the
\texttt{timm} backbone in this mode (\osnetGT{} vs.\ \timmGT{} HOTA) isolates how much association
quality the appearance model contributes once detection is held perfect; comparing the same models in
the full pipeline against this ceiling reveals how detector-bound association quality is in practice.

\paragraph{(b) Per-video detector-confidence sweep.}
For each of the six videos we sweep the YOLO confidence threshold over
$\{0.15, 0.25, 0.35, 0.45\}$, holding the detector, REID model, input size, and SORT core fixed. This
is the precision/recall lever: lowering the threshold raises recall and lowers precision. We pair the
HOTA response with the detector's per-clip precision at conf~0.25 to test whether the \emph{magnitude}
of recall's effect on HOTA is predicted by per-clip precision.

\paragraph{(c) Input-size sweep on the two dense clips.}
On the two dense, low-resolution clips (TUD-Campus and KITTI-17), where YOLO precision is low, we
additionally sweep the detector input resolution over $\{960, 1280, 1536\}$ (the $1280$ midpoint is the
default configuration, shared with the confidence sweep). Larger input recovers
more small instances -- a second, \emph{orthogonal} recall lever to the confidence threshold -- and
lets us test whether it reproduces the confidence-sweep effect or, by changing \emph{which} detections
exist rather than merely gating them, behaves differently.

\paragraph{(d) Gating-restoration ablation.}
To attribute the regression specifically to the dropped gate, we restore a DeepSORT-style admission
gate -- the reference reproduction value $\text{min\_confidence}=0.3$~\cite{wojke2017deepsortcode}
together with secondary NMS at overlap~$0.7$ (our setting; the reference default is $1.0$) -- on top of
the otherwise unchanged modern pipeline, and re-measure HOTA on each video. Because the gate acts only
on the detection stream feeding the unchanged tracker, any recovered HOTA is causally attributable to
admission, not to a confounding code change.

\paragraph{Reproducibility.}
The detector is pinned: ultralytics~8.4.79 with the \texttt{yolov8m.pt} checkpoint
(sha256 \texttt{5d4a90cd\dots4fe5}) at input size~1280. Per-clip detector precision is
weight-version-sensitive, so we fix and report the version; a different ultralytics release
shifts the absolute values (e.g.\ TUD-Campus precision $0.50$ rather than $0.43$) but leaves the
low-precision ordering -- TUD-Campus and KITTI-17 lowest -- on which the analysis depends unchanged.
Every run -- each detector preset, REID model, confidence level, input size, and ablation -- appends
one row to a versioned CSV journal recording the configuration and resulting metrics, so the full
experimental history is traceable and the gating-sweep figures in \cref{sec:results} (confidence,
precision, input-size, and ablation plots) are regenerable from it; the decomposition and
ground-truth-box numbers come from separate evaluation runs recorded in the same journal. The
sweep harness writes each result immediately and skips configurations already present, so an
interrupted run resumes without recomputation. We also note two environment-level pitfalls that the
journal helped surface and that we report as negative results in \cref{sec:limitations}: a NumPy-2.x
incompatibility that breaks the torchreid/boxmot/TrackEval build chain, and a 2D~MOT~2015-vs-MOT16
ground-truth column-format mismatch that silently zeroed four sequences until corrected.
